% !TEX root = ../main.tex
\chapter{Testing}

The definition of done set out in Chapter 2 required every task to be covered by automated tests, so the suite grew
with the code rather than after it. The standard suite currently counts \textbf{515 tests over 61 suites}. Each
member wrote the tests of the epics or the tasks they were responsible for, so the ownership of an area covered its
specification as well as its implementation.

\section{Conventions}

As a rule the tests of a task were written before its implementation, and the history of the commits keeps the trace of
it.

Suites are written in the \texttt{AnyFunSuite} style, which lets the name of a test be a sentence rather than an
identifier. The convention adopted is that a name describes \textbf{a behaviour of the domain} and not the method under
test, so that it stays true when the method is renamed. Each test covers \textbf{one concept} and follows the
build-operate-check shape.

What several suites share lives in six dedicated files that hold no test of their own: a maze for testing together with
the levels derived from it (\texttt{LevelTestSupport}), the mazes drawn in ASCII the strategies are tried on
(\texttt{MazeLayouts}), the reading of test mazes from the resources (\texttt{MapTestSupport}), sample results and
leaderboards at a fixed instant (\texttt{LeaderboardTestSupport}), a fake state with a listener recording what it is
notified of (\texttt{RenderingTestSupport}) and the checker of the monoid laws (\texttt{MonoidLawsTest}).

\section{Testing the pure model}

The rules of the game are pure functions from a state and a duration to the state that follows, so they are tested
without infrastructure: there is nothing to set up and nothing to tear down.

\subsection{Loop and update pipeline}

The clock is verified as an immutable component that accumulates elapsed time and rejects negative advances. The loop
lifecycle is tested as a state machine, covering the valid transitions and the invalid ones, which must raise an
exception. The update pipeline is covered with order-focused tests ensuring that each tick runs the expected stages in
sequence and that each stage receives the result of the previous one. The outcome of a level is modelled apart from the
lifecycle of the loop and is tested for the property the rest of the code relies upon: victory and defeat are terminal,
running is not.

\subsection{Maps}

Map validation is one of the larger suites of the project: each rule listed under FR-S2 has a test that rejects a maze
breaking it, and the errors are asserted to \textbf{accumulate} rather than stopping at the first, since a player
rewriting a maze should be told everything that is wrong with it at once. The mazes under test are resources on file
rather than strings in the code, which is what makes it practical to cover every documented teleport pairing.

\subsection{Movement and collisions}

Movement is tested as an entity crossing between adjacent cells over a given duration: the crossing in progress, its
completion, the direction faced independently of the movement under way, and the rejection of a negative duration or of
a non-positive time per cell. Collision detection is tested by cell, covering every kind of occupant and the case of a
tile and an enemy on the same position. Two entities exchanging cells within one tick are covered as a meeting, which
is the case a position-by-position comparison alone would miss.

\subsection{Enemy behaviour}

The two pursuing strategies are tested on hand-built mazes: they take the shortest path around dead ends, respect walls
when the closest move is blocked, return no move when trapped, and use a teleport only when it shortens the path.
Determinism is asserted explicitly, so that two equally close moves always resolve the same way, and one test asserts
that pursuit and anticipation produce observably different decisions, which is what makes them worth keeping apart.

The patrolling strategy is tested for a behaviour no single step can show. One test walks a simulated patrol for sixty
moves through a tail-recursive traversal and compares the corners touched with the endless repetition of the expected
round, establishing that the route closes and repeats. Another asserts that its choices do not change when the player
is moved elsewhere.

\subsection{Rules and effects}

Every rule depending on elapsed time receives it as an argument instead of reading a clock, which makes temporal
behaviour testable \textbf{without waiting}: a bonus is asserted to be still applied one millisecond before its
expiration and no longer applied at it, pinning the exact boundary. The same technique covers effects renewed while
already active, effects expiring independently of one another, and the restoration of an effect at a different instant
of the clock, which is what a saved game needs when it is taken up again.

Where a rule is a property rather than an example, the test states the property: collection is idempotent, so an
emptied position yields nothing and leaves the count of what remains untouched; a maze holding only bonuses is complete
from the start. Meeting an enemy is covered on both sides of the rule, and the slowdown effect is asserted not to grant
the protection that invulnerability does.

\subsection{Modes and scoring}

Each mode is tested for the three things that distinguish it: whether it has a clock and how much of it is left, what
decides its outcome, and how it awards points. Survival is covered for its progressive increase in enemy speed and for
the order in which that increase and the slowdown effect combine. Scoring is tested through tables of event and
expected points, including the combo that is reset without affecting the points already gained; each scoring rule is
asserted to award nothing for the events belonging to the others.

The leaderboard is tested as an algebraic structure: a reusable checker verifies the \textbf{monoid laws}, left and
right identity and associativity, over a list of leaderboards including the empty one and one exceeding the cap. Its
ordering is covered for scores and, at equal scores, for the instant they were achieved.

\subsection{Persistence}

Saved games and recorded results are tested for round trips and for failures. A saved game is written and read back
with all its progress, down to the movement an enemy was making and the corner a patroller was heading for. Failures
are covered by \textbf{injecting them}: a real save is written, one property of the file is then corrupted, and the
load is required to refuse it without throwing, for an unsupported format version, an invalid mode, an unsafe maze
name, a negative number of lives and an entity outside the map. Recorded results are covered by the same round trip on
their encoding, \texttt{decode(encode(r)) == r}, together with the malformed lines that must be rejected.

\section{Testing the interface without opening a window}

Continuous integration runs on a headless machine, so a test instantiating a graphical node would break the build for
the whole team. This is why the interface is split as Chapter 5 describes, and everything on the pure side of that
split is tested: which screen is showing, what covers the board, how the maze is laid out, how wide a cell is drawn,
what the status bar reads, what the standings are.

The state of the level is passed to the overlay by name, because a game in progress never needs it, and a test counts
the evaluations of that argument to assert that a board covered by nothing does not even ask how the level is doing.
Other tests assert that every sprite the game can name has a picture among the resources and that every teleport pair a
maze may hold is one the board knows how to draw, catching at build time a kind of mismatch that usually appears only
when the application is run.

\section{Testing the boundary}

The interface through which the game reaches the world outside is the seam at which that world is replaced during
testing. The tests of the application supply an in-memory implementation of it, which answers what it is told to
answer, remembers what it was asked to keep, and can be instructed to fail on reading or on writing. Behaviours that
would otherwise require a real game and a real file system are then asserted directly: a result is recorded
\textbf{once}, however many frames follow the end of a game; a game whose save failed is not put away and its player is
told; a game resumed from a file has no leaderboard to be recorded in; a maze opened from elsewhere is kept, so that it
is offered from then on. The controller is thus tested end to end \textbf{without a file being written and without a
window being opened}. The adapters that do touch the file system are tested separately, against temporary directories,
so that the decision and the storage fail apart from each other.

\section{Levels of testing}

\begin{itemize}
  \item \textbf{Unit}: the large majority of the suite, each test covering one abstraction at a time.
  \item \textbf{Integration}: a single test, marked \textbf{by convention} with the \texttt{IT} suffix, exercising
  movement over a validated maze rather than over a grid built for the occasion, so that two modules are put to the
  test together. It runs with all the others.
  \item \textbf{Performance}: the tick check described among the non-functional requirements, which lives on sources of
  its own under a dedicated sbt configuration and is therefore left out of the ordinary suite. It is run separately, and
  is what measures the margin against the sixteen milliseconds a frame is given.
  \item \textbf{Acceptance}: not automated, carried out by the member acting as client and domain expert, who played
  each increment at the end of its sprint and reported what the specification did not cover or what the game did not do
  as expected.
\end{itemize}

\section{Design decisions the tests pin down}

Two values of the implementation are written down in the tests alone: the code applies them without explaining them,
and changing one turns a test red instead of silently changing how the game behaves.

Renderings notify their listeners only after a \textbf{visible} change, so that a frame in which nothing changed is not
redrawn. For that comparison to be worth anything, the elapsed time the projection exposes is \textbf{truncated to whole
seconds}: whole seconds are what the status bar displays, and an untruncated time would differ at every tick, leaving
nothing to discard. In the code these are two barely visible calls, the time elapsed rounded down and the time left
rounded up, and what says why they are there is the test requiring that a tick with nothing visible to show notifies
nobody.

The clamp on the duration of a frame is the same case. At every frame the model receives the time passed since the
previous one, ordinarily sixteen milliseconds, but that measure is not to be trusted: it can arrive negative when a
frame presents itself out of order, or amount to whole seconds when the machine has stalled. The clock of the level is
what the expiry of the effects and the time left in a timed game rest upon, so a long frame would expire a bonus just
picked up while nothing happened on the screen. The time handed to the model is therefore kept between zero and fifty
milliseconds, and the tests fix both limits: the clock does not go back for a frame out of order, and does not advance
by more than one step for a frame arriving late.

\section{Coverage of the requirements}

Every functional requirement of Chapter 3 is answered by tests. The table names, for each one, the suites that establish
it; the non-functional requirements name their own evidence where they are stated.

\begin{center}
  \small
  \begin{tabular}{|p{1.6cm}|p{4.4cm}|p{6.4cm}|}
    \hline
    \textbf{Req.} & \textbf{Requirement} & \textbf{Where it is established} \\
    \hline
    FR-U1 & Give a name & \texttt{GameFilesTest}, \texttt{MenuScreenTest} \\
    \hline
    FR-U2 & Choose a maze & \texttt{PlayableMazesTest}, \texttt{DefaultMapsTest}, \texttt{ApplicationTest} \\
    \hline
    FR-U3 & Choose a mode & \texttt{ModeTuningTest}, \texttt{ApplicationTest} \\
    \hline
    FR-U4 & Start a game & \texttt{ApplicationTest} \\
    \hline
    FR-U5 & Move through the maze & \texttt{LevelMovementTest}, \texttt{CommandMapperTest} \\
    \hline
    FR-U6 & Follow the game as it is played & \texttt{StatusBarTest}, \texttt{LevelViewTest}, \texttt{FrameTest} \\
    \hline
    FR-U7 & Pause, resume and load & \texttt{GameSessionTest}, \texttt{ApplicationTest},
    \texttt{PropertiesGameSaveRepositorySpec} \\
    \hline
    FR-U8 & Restart & \texttt{ApplicationTest} \\
    \hline
    FR-U9 & See the leaderboard & \texttt{StandingsTest}, \texttt{LeaderboardTest} \\
    \hline
    FR-S1 & Read a maze & \texttt{MapParserSpec}, \texttt{MapLoaderSpec} \\
    \hline
    FR-S2 & Judge whether a maze is playable & \texttt{MapValidationSpec}, \texttt{MapPipelineSpec} \\
    \hline
    FR-S3 & Explain a refusal & \texttt{MapValidationSpec}, \texttt{GameFilesEnvironmentTest} \\
    \hline
    FR-S4 & Movement & \texttt{MovingEntityTest}, \texttt{LevelMovementTest} \\
    \hline
    FR-S5 & Decide where enemies go & \texttt{EnemyMovementStrategySpec}, \texttt{EnemyMovementSpec},
    \texttt{PatrolStrategySpec}, \texttt{PatrolRouteSpec}, \texttt{EnemyAiStageSpec},
    \texttt{EnemyStrategySelectorSpec} \\
    \hline
    FR-S6 & Carry entities through teleports & \texttt{LevelStateTest}, \texttt{MovingEntityMapIT} \\
    \hline
    FR-S7 & Resolve a meeting & \texttt{LevelMeetingTest}, \texttt{CollidingTest}, \texttt{CollisionDetectorTest},
    \texttt{CollisionResolverTest} \\
    \hline
    FR-S8 & Grant and expire effects & \texttt{ActiveEffectsTest}, \texttt{BonusCollectingTest} \\
    \hline
    FR-S9 & Award points & \texttt{LevelScoreTest}, \texttt{ScoreTrackerTest}, \texttt{ScoringRuleTest} \\
    \hline
    FR-S10 & Judge game status & \texttt{LevelOutcomeTest}, \texttt{GameModeSpec} \\
    \hline
    FR-S11 & Keep the best results only & \texttt{LeaderboardTest}, \texttt{FileLeaderboardStorageTest},
    \texttt{GameResultCodecTest}, \texttt{LeaderboardRecordingTest} \\
    \hline
  \end{tabular}
\end{center}
